\subsection{Representations and the Grothendieck Group}

In this subsection, we assume that K is a commutative field. Given a finite-dimensional linear representation \((M,\pi)\) of G, we denote by \([\pi]\) the class of the K{[}G{]}-module M in the Grothendieck ring \(\mathrm{R}_{\mathrm{K}}(\mathrm{G})\) (VIII, p.~198, Example 1).

By the definition of the laws of composition on the ring \(\mathrm{R}_{\mathrm{K}}(\mathrm{G})\) , if \(\pi\) and \(\pi'\) are finite-dimensional linear representations of G, then we have

\[
[ \pi \oplus \pi^ {\prime} ] = [ \pi ] + [ \pi^ {\prime} ], \quad [ \pi \otimes \pi^ {\prime} ] = [ \pi ] [ \pi^ {\prime} ].
\]

The unit element of the ring \(R_{K}(G)\) is the class of the unit representation of G (VIII, p.~399, Example 1).

The ring \(R_{K}(G)\) is a free Z-module with basis the set of classes of simple K{[}G{]}-modules of finite dimension over K. When \(|G|\) is invertible in K, two finite-dimensional representations are isomorphic if and only if they have the same class in \(R_{K}(G)\) (VIII, p.~190, Corollary and p.~401, Corollary 1).

For every finite-dimensional linear representation \((M,\pi)\) of G, the contragredient \(\pi^{\vee}\) is also finite-dimensional. The representations \(\pi\) and \((\pi^{\vee})^{\vee}\) are isomorphic. If \(\pi'\) is also a finite-dimensional linear representation of G, then the representations \((\pi\otimes\pi')^{\vee}\) and \(\pi^{\vee}\otimes\pi'^{\vee}\) are isomorphic. If

\[
0 \longrightarrow \mathrm{M} ^ {\prime} \longrightarrow \mathrm{M} \longrightarrow \mathrm{M} ^ {\prime \prime} \longrightarrow 0
\]

is an exact sequence of K{[}G{]}-modules of finite dimension over K, then the contragredient representations give an exact sequence of K{[}G{]}-modules

\[
0 \longrightarrow \mathrm{M} ^ {\prime \prime \vee} \longrightarrow \mathrm{M} ^ {\vee} \longrightarrow \mathrm{M} ^ {\prime \vee} \longrightarrow 0
\]

(II, §7, No.~5, p.~299, Corollary of Theorem 5). By the universal property of Grothendieck groups (VIII, p.~186, Proposition 4), there exists an automorphism \(c \mapsto c^{\vee}\) of the ring \(\mathrm{R}_{\mathrm{K}}(\mathrm{G})\) characterized by \([\pi]^{\vee} = [\pi^{\vee}]\) for every finite-dimensional linear representation \(\pi\) of \(\mathrm{G}\) ; we have \((c^{\vee})^{\vee} = c\) for every element \(c\) of \(\mathrm{R}_{\mathrm{K}}(\mathrm{G})\) .

Let H be a subgroup of G. Restriction preserves exact sequences. By the universal property of Grothendieck groups (loc. cit.), there consequently exists a group homomorphism, denoted by \(Res_{H}^{G}\) , from \(R_{K}(G)\) to \(R_{K}(H)\) characterized by the relation

\[
\mathrm{Res} _ {\mathrm{H}} ^ {\mathrm{G}} [ \pi ] = [ \mathrm{Res} _ {\mathrm{H}} ^ {\mathrm{G}} (\pi) ]\tag{6}
\]

for every finite-dimensional linear representation \(\pi\) of G; it is a ring homomorphism. If H has finite index in G, then by Proposition 14 of II, §7, No.~7, p.~306, we can define, analogously, a group homomorphism \(\mathrm{Ind}_{\mathrm{H}}^{\mathrm{G}}\) from \(\mathrm{R}_{\mathrm{K}}(\mathrm{H})\) to \(\mathrm{R}_{\mathrm{K}}(\mathrm{G})\) characterized by the relation

\[
\mathrm{Ind} _ {\mathrm{H}} ^ {\mathrm{G}} [ \sigma ] = [ \mathrm{Ind} _ {\mathrm{H}} ^ {\mathrm{G}} (\sigma) ]\tag{7}
\]

for every finite-dimensional linear representation \(\sigma\) of H.

By relations (1) of VIII, p.~399 and (3) of VIII, p.~400 and the universal property of Grothendieck groups mentioned above, there exists a ring homomorphism \(\Theta_{\mathrm{G}}\) from \(\mathrm{R}_{\mathrm{K}}(\mathrm{G})\) to the algebra \(\mathcal{Z}_{\mathrm{K}}(\mathrm{G})\) of central functions on \(\mathrm{G}\) , characterized by \(\Theta_{\mathrm{G}}([\pi]) = \chi_{\pi}\) for every finite-dimensional representation \(\pi\) of \(\mathrm{G}\) .

If H is a subgroup of G, then the homomorphisms \(\Theta_{G}\) and \(\Theta_{H}\) associated with the groups G and H are compatible with the operations \(Res_{H}^{G}\) and \(Ind_{H}^{G}\) (VIII, p.~402 and VIII, p.~404, Proposition 3).

Suppose that G is the product \(G^{\prime} \times G^{\prime\prime}\) of two groups. By VIII, p.~213, Remark 2, there exists a Z-linear mapping \(\kappa\) from \(\mathrm{R}_{\mathrm{K}}(\mathrm{G}^{\prime}) \otimes_{\mathbf{Z}} \mathrm{R}_{\mathrm{K}}(\mathrm{G}^{\prime\prime})\) to the group \(\mathrm{R}_{\mathrm{K}}(\mathrm{G}^{\prime} \times \mathrm{G}^{\prime\prime})\) characterized by the relation \(\kappa([\pi^{\prime}] \otimes [\pi^{\prime\prime}]) = [\pi^{\prime} \boxtimes \pi^{\prime\prime}]\) for \(\pi^{\prime}\) (resp. \(\pi^{\prime\prime}\) ) a finite-dimensional representation of \(G^{\prime}\) (resp. \(G^{\prime\prime}\) ) and \(\pi^{\prime} \boxtimes \pi^{\prime\prime}\) the external tensor product (VIII, p.~400, Example 5). It is a ring homomorphism. If the field K is algebraically closed, then the mapping \(\kappa\) is an isomorphism (VIII, p.~213, Remark 2).

Suppose that G is the product \(G' \times G''\) of two groups. Denote by \(\psi\) the homomorphism from \(\mathcal{Z}_{\mathrm{K}}(\mathrm{G}') \otimes_{\mathrm{K}} \mathcal{Z}_{\mathrm{K}}(\mathrm{G}'' )\) to \(\mathcal{Z}_{\mathrm{K}}(\mathrm{G})\) that sends \(f' \otimes f''\) to the function \((g', g'') \mapsto f'(g')f''(g'')\) . The following diagram commutes:

\[
\mathrm{R} _ {\mathrm{K}} (\mathrm{G} ^ {\prime}) \otimes_ {\mathbf {Z}} \mathrm{R} _ {\mathrm{K}} (\mathrm{G} ^ {\prime \prime}) \xrightarrow {\kappa} \mathrm{R} _ {\mathrm{K}} (\mathrm{G})
\]

\[
\bigg \downarrow \Theta_ {\mathrm{G} ^ {\prime}} \otimes \Theta_ {\mathrm{G} ^ {\prime \prime}} \qquad \qquad \bigg \downarrow \Theta_ {\mathrm{G}}
\]

\[
\mathcal {Z} _ {\mathrm{K}} (\mathrm{G} ^ {\prime}) \otimes_ {\mathrm{K}} \mathcal {Z} _ {\mathrm{K}} (\mathrm{G} ^ {\prime \prime}) \xrightarrow {\psi} \mathcal {Z} _ {\mathrm{K}} (\mathrm{G}).
\]

